\section{Discusión}
\label{sec:discusion}

\subsection*{RQ1: ¿Puede automatizarse el flujo completo con trazabilidad verificable?}
Sí, parcialmente verificado. Los 90 requisitos (64 funcionales, 26 no funcionales; SRS v1.0.1) están
implementados o planificados según su estado declarado, y el flujo completo
(solicitud $\to$ anteproyecto $\to$ jurado $\to$ cronograma $\to$ evaluación $\to$
acta) se verificó manualmente end-to-end contra la topología de producción. La trazabilidad
\textbf{automatizada} cubre 36 de 50 requisitos Must con prueba dedicada (72\,\%; Sección~\ref{sec:requisitos});
los 14 Must restantes se enumeran uno por uno, con lo que falta en cada uno, en la matriz de trazabilidad
(\texttt{docs/trazabilidad/matriz.csv}, Anexo~\ref{anexo:matriz}). La respuesta a RQ1 es ``sí, funcionalmente'' y
``72\,\%, en términos de verificación automatizada de los requisitos Must''.

\subsection*{RQ2: ¿Cuál es el efecto cuantificable de la caché sobre la latencia?}
Respondida con evidencia estadística fuerte: reducción de latencia de 13.9$\times$ (109.62\,ms a
7.86\,ms), estadísticamente significativa ($p = 3.02\times10^{-11}$) con separación perfecta entre
distribuciones (tamaño de efecto $r=-1.00$, magnitud máxima). Es la pregunta de investigación con la respuesta más sólida
de las cuatro, porque es la única con un diseño experimental que aísla una sola variable (estado de la
caché) manteniendo todo lo demás constante.

\subsection*{RQ3: ¿Qué vulnerabilidades reales existen y qué tan cerradas están?}
Ninguna vulnerabilidad de inyección SQL (0 hallazgos de \texttt{SQL\_NONCONSTANT\_STRING\_PASSED\_TO\_EXECUTE}
en 233 hallazgos totales de find-sec-bugs al 29-08), 0 fallos \texttt{FAIL} y 0 alertas de riesgo \texttt{High} en
el escaneo dinámico OWASP ZAP (la única \texttt{High} encontrada, \texttt{@angular/core} vulnerable, se
corrigió el 2026-08-29). Sin embargo, existen 14 vulnerabilidades de dependencias de terceros sin
corregir (5 críticas, todas en herramientas de build de íconos que no llegan a producción) y 9 hallazgos
de \texttt{PATH\_TRAVERSAL\_IN} sin resolver (riesgo real bajo, porque las entradas son IDs \texttt{Long},
no cadenas arbitrarias, pero no verificado formalmente).

Ese panorama, sin embargo, solo cubre lo que find-sec-bugs y ZAP pueden detectar por patrón sintáctico
--- ninguna de las dos herramientas evalúa \emph{lógica} de autorización por recurso. Una revisión manual
dirigida encontró y corrigió 6 fallos reales de ese tipo que las herramientas automatizadas no habían
señalado: IDOR de lectura en evaluaciones y anteproyectos, IDOR de escritura en el registro de notas de
jurado, IDOR en el perfil de docentes, un endpoint de eliminación de actas que exigía solo estar
autenticado en vez de un permiso administrativo, mass-assignment en el registro/creación de usuarios que
permitía sobrescribir una cuenta existente, y secuencias de PostgreSQL desincronizadas (detalle completo
en Sección~\ref{sec:seguridad-corregida}). Los 6 están corregidos, verificados en vivo contra el backend
real, y cubiertos con pruebas automatizadas nuevas. La respuesta honesta a RQ3 es entonces de dos capas:
las vulnerabilidades de \textbf{código propio y de código servido al navegador} detectables por
herramienta automatizada (inyección SQL, comparación de hash insegura, librería Angular vulnerable) están
cerradas y verificadas; las de \textbf{lógica de autorización}, invisibles a esas herramientas, también se
buscaron activamente y las 6 encontradas se cerraron; las vulnerabilidades restantes conocidas están en
\textbf{herramientas de desarrollo}, no en código que se ejecuta en producción.

\subsection*{RQ4: ¿Qué tan honesta puede mantenerse la documentación bajo presión de plazos?}
Esta pregunta se responde con el propio historial del proyecto, no con una medición externa. Se
detectaron y corrigieron, dentro de este mismo proceso: una encuesta SUS fabricada (con evaluadores
inexistentes, ver \texttt{docs/mediciones/sus/SUS-RESULTS.md}), métricas Lighthouse citadas sin haber corrido la herramienta, 5 citas de
archivos de test inexistentes en la matriz de trazabilidad, una afirmación falsa sobre anclaje de
imágenes Docker por digest, y una afirmación falsa sobre custodia de formularios de consentimiento
firmados. En cada caso, la corrección fue reemplazar el dato fabricado por una medición real o una
declaración explícita de ausencia --- nunca por otro dato fabricado más plausible. Esto sugiere que sí es
posible mantener honestidad bajo presión de plazos, al costo de reportar cifras menos favorables en su
momento (63.17\,\% de cobertura de líneas en el cierre del 2026-09-05, hoy 82.01\,\% global y 85.1\,\% de ramas en controladores, en
vez de un ``$>$60\,\%'' fabricado; SUS ``pendiente'' en vez de un puntaje fabricado, hoy aplicado con
$n=15$ con fecha sellada por un tercero en vez de reportar las 15 hojas recolectadas cuando 11 tienen una fecha que no se
sostiene --- ver \texttt{docs/mediciones/sus/SUS-RESULTS.md}), lo cual es precisamente el trade-off que
esta guía exige explícitamente priorizar.
\newpage
